\documentclass{article}
\usepackage{tcee-repaso}
\tcetema[Modelización dinámica]{3.A.29}{Modelización dinámica de las tomas de decisiones. Modelos de horizonte infinito y modelos de generaciones solapadas.}
\tcefecha{03/11/2026}
\begin{document}
\cabecerarepaso
\begin{columnas}

% ==========================================================================
\bloque{Esquema del cante}
\esquemacante{%
\apartadocante{Introducción: «no dejes para mañana\ldots»}{2,5}
\apartadocante{I. La modelización dinámica: enfoque micro}{8,5}
\subcante{1. Consumidores: componentes, 3 métodos, \autor{Fisher}}{6}
\subcante{2. Empresas: \autor{Fisher}, q de \autor{Tobin}, \autor{Hotelling}}{2,5}
\apartadocante{II. El agente representativo de vida infinita}{8}
\subcante{1. El agente representativo}{2}
\subcante{2. Horizonte infinito y \autor{Ramsey}-\autor{Cass}-\autor{Koopmans}}{6}
\apartadocante{III. Generaciones solapadas y heterogeneidad}{10}
\subcante{1. \autor{Samuelson} (1958) y \autor{Diamond} (1965)}{8}
\subcante{2. Limitaciones y nuevos enfoques}{2}
\apartadocante{Conclusión}{1}}
\emph{Pizarra} en 3 columnas: izquierda, recuadrados hasta el final, el problema del consumidor, \autor{Euler}, \autor{Bellman} y la ley de movimiento de \autor{Diamond}; centro, tabla de los 3 métodos (se borra al acabar I) y luego las 3 ecuaciones de RCK; derecha, gráfico de \autor{Diamond} y cuadro MHI vs. MGS. Pegatinas: «cada cosa a su tiempo» (\autor{Bellman}), «de tal palo, tal astilla» (dinastía), «cuando seas padre, comerás huevos» (MGS), «la avaricia rompe el saco» (ineficiencia dinámica).

% ==========================================================================
\bloque[2,5]{Introducción}
\begin{itemize}
\item \textbf{Enganche}: «No dejes para mañana lo que puedas hacer hoy»; la economía dice «depende»: si lo ahorrado rinde más de lo que cuesta esperar, conviene esperar. Y el lunes de empezar a ahorrar pasa al siguiente. \autor{Ramsey} (1928, 25 años): ¿cuánto de su renta debe ahorrar una nación?
\item \textbf{Relevancia}: casi todas las decisiones importantes son dinámicas (ahorrar, invertir, educarse, endeudarse, extraer, cotizar); lenguaje de la macro actual desde \autor{Lucas}; reparto entre generaciones: \concepto{mecanismo de equidad intergeneracional} (desde 2023), 0,90 p.p. en 2026 $\to$ 1,2 en 2029, hasta 2050, al Fondo de Reserva, sin derecho a prestación.
\item \textbf{Contextualización}: \concepto{análisis estático} (variables de un mismo momento) vs. \concepto{dinámico} (de momentos distintos: subíndices temporales); \concepto{estática comparativa} entre ambos; \concepto{principio de correspondencia} (\autor{Samuelson}, 1947): la estática comparativa solo tiene sentido si el equilibrio es estable. \autor{Tinbergen}: 2 formas de estudiar, no 2 fenómenos; aporta si la decisión es \concepto{genuinamente intertemporal}. Historia: clásicos mezclan; marginalistas, estática; \autor{Fisher} (1930), \autor{Ramsey} (1928), \autor{Frisch} (1933); microfundamentos (\autor{Friedman}, \autor{Modigliani}, \autor{Jorgenson}, \autor{Baumol}-\autor{Tobin}, \autor{Lucas}-\autor{Rapping}); \autor{Bellman} (1957), \autor{Pontryagin} et al. (1962); MHI (\autor{Cass}, \autor{Koopmans}, 1965) y MGS (\autor{Samuelson}, 1958; \autor{Diamond}, 1965); tras \autor{Lucas} (1976), agentes representativos y luego heterogéneos.
\item \textbf{Problemática}: ¿cómo deciden cuando hoy cambia mañana, y con qué herramienta? ¿Qué cambia entre un agente inmortal y generaciones que se suceden? ¿Ahorra el mercado lo que conviene?
\item \textbf{Estructura}: de lo individual a lo agregado: \hyperref[rep:I]{I}, la herramienta (3 métodos, una \autor{Euler}); \hyperref[rep:II]{II}, el agente de vida infinita (RCK, eficiente); \hyperref[rep:III]{III}, generaciones solapadas (puede ahorrar de más) y límites. Fuera: juegos dinámicos \vertema{3.A.14}.
\end{itemize}

% ==========================================================================
\bloque[8,5]{I. La modelización dinámica: enfoque micro\label{rep:I}}
\begin{ideaclave}
Un problema es genuinamente dinámico cuando una variable de estado une los periodos. Con cualquiera de los 3 métodos, la solución es una condición de arbitraje, la ecuación de \autor{Euler}, y un precio sombra del estado.
\end{ideaclave}

\sub{I.1. Decisiones intertemporales de los consumidores\label{rep:I.1}}
\begin{itemize}
\item \concepto{Variables flujo} = \concepto{de control} ($C_t$, $I_t$); \concepto{variables stock} = \concepto{de estado} ($b_t$, $K_t$, reservas). 4 piezas: preferencias sobre sendas, restricciones, método, senda óptima.
\item \emph{Preferencias}: $U_0=U(C_0,C_1,\dots)$; \autor{Ramsey}: $\sum D(t)u(C_t)$; \autor{Samuelson} (1937), \concepto{descuento exponencial}:
\begin{formula}
\centering $U_0\equiv\sum_{t=0}^{T}\beta^t u(C_t),\quad \beta\equiv\dfrac1{1+\rho},\ \rho>0,\ 0<\beta<1$
\end{formula}
\item[] $u'>0$, $u''<0$ (alisar), bienes no duraderos. \concepto{ESI} $\equiv d\ln(C_{t+1}/C_t)/d\ln|\mathrm{RMS}_{t,t+1}|$. \concepto{CRRA}:
\begin{formula}
\centering $u(C)=\dfrac{C^{1-\theta}-1}{1-\theta}\ (\theta>0,\theta\neq1);\ \ \ln C\ (\theta=1);\ \ \mathrm{ESI}=\dfrac1\theta$
\end{formula}
\item[] $\theta$ = aversión relativa al riesgo: misma aversión a variabilidad entre estados y entre periodos (\autor{Epstein}-\autor{Zin} las separan).
\item \emph{Propiedades} (MWG, 20.B): impaciencia; estacionariedad; separabilidad aditiva (sin hábitos); duración del periodo; recursividad $U_0=u(C_0)+\beta U_1$; altruismo (dinastía). Impaciencia + estacionariedad $\Rightarrow$ exponencial, el único consistente. Lo intertemporal viene de la \emph{restricción}.
\item \emph{Restricción} por periodo (ecuación de estado):
\begin{formula}
\centering $C_t+b_{t+1}=Y_t+(1+r)b_t\iff b_{t+1}-b_t=Y_t+rb_t-C_t$
\end{formula}
\item[] $b_0$ dado, $b_{T+1}\geq0$ (= 0 en el óptimo); con $T=\infty$, \emph{no-Ponzi}. Intertemporal: VA del consumo = $(1+r)b_0$ + VA de las rentas (con incertidumbre, exige mercados completos; activos redundantes: indeterminada la cartera, no el consumo).
\item \textbf{a) \autor{Lagrange}} (discreto, finito): CPO $\beta^tu'(C_t)=\lambda/(1+r)^t$; cociente $\Rightarrow$ \concepto{ecuación de \autor{Euler}}:
\funcionclave{u'(C_t)=\beta\,(1+r)\,u'(C_{t+1})}
\item[] Utilidad perdida hoy = la que dan mañana el euro y sus intereses, descontada. $\lambda$ = utilidad marginal de la riqueza. Límite: muchas condiciones e incertidumbre.
\item \textbf{b) Control óptimo} (continuo; \concepto{principio del máximo}: \autor{Boltianski}-\autor{Gamkrelidze}-\autor{Pontriaguin}, 1956; publicado 1962). $\mathcal H_t=u(C_t)+\lambda_t(Y_t+rb_t-C_t)$: $u'(C_t)=\lambda_t$; $\dot\lambda_t=(\rho-r)\lambda_t$; transversalidad $\lim e^{-\rho t}\lambda_tb_t=0$. Con CRRA:
\begin{formula}
\centering $\dfrac{\dot C_t}{C_t}=\dfrac{r-\rho}{\theta}$
\end{formula}
\item[] Ventaja: coestado = \concepto{precio sombra} (q, regalía); diagramas de fases (RCK, inversión, empleo).
\item \textbf{c) Programación dinámica} (\autor{Bellman}, 1957): \concepto{principio de optimalidad}: el resto de un plan óptimo es óptimo desde el estado al que lleva: «cada cosa a su tiempo». \concepto{Función de valor} $V(b_t)$ (indirecta, como $V(\vect p,\overline W)$):
\funcionclave{V(b_t)=\max_{C_t}\{u(C_t)+\beta V(b_{t+1})\}}
\item[] CPO $u'(C_t)=\beta V'(b_{t+1})$ + envolvente $V'(b_t)=(1+r)u'(C_t)$ $\Rightarrow$ \autor{Euler}. $V'$ = $\lambda$ = coestado.
\item[] \emph{Ejemplo}: $Y_t=0$, $\ln C$, $T=\infty$; se prueba $V=A+B\ln b$: $B=1/(1-\beta)$,
\begin{formula}
\centering $C_t=(1-\beta)(1+r)\,b_t,\qquad b_{t+1}=\beta(1+r)\,b_t$
\end{formula}
\item[] Fracción fija $1-\beta$, independiente de $r$ (con $\ln$, renta y sustitución se compensan). Ventajas: incertidumbre, cálculo numérico (converge por $\beta<1$), regla $C(b_t)$: ciclo real, EGDE, búsqueda, \autor{Shapiro}-\autor{Stiglitz}.
\item \textbf{d) Incertidumbre}: $V(b_t,Y_t)=\max\{u+\beta\mathbb E_tV(\cdot)\}$ y $u'(C_t)=\beta(1+r)\mathbb E_t[u'(C_{t+1})]$; una restricción por periodo y estado; expectativas racionales.
\item \textbf{e) Tabla}: \emph{Lagrange}: discreto, finito; senda entera; $\lambda$; sencillez; \autor{Fisher}, ciclo vital, \autor{Diamond}. \emph{Control}: continuo; senda del control; $\mathcal H$, $\lambda_t$; fases; q, RCK, empleo. \emph{PD}: discreto; regla $C(b)$; $V(b)$; incertidumbre; \autor{Hall}, ciclo real, búsqueda. Mismo resultado: \autor{Euler}.
\item \textbf{\autor{Fisher}} (1930), 2 periodos, $C_1+C_2/(1+r)=\overline W$; con CRRA:
\begin{formula}
\centering $\dfrac{C_{t+1}}{C_t}=[\beta(1+r)]^{1/\theta}\ \Rightarrow\ \ln\dfrac{C_{t+1}}{C_t}\simeq\dfrac{r-\rho}{\theta}$
\end{formula}
\item[] $r\gtrless\rho$: el consumo crece/decrece; no dice si presta (depende de la senda de la renta). ESI baja (\autor{Hall}, 1988; \autor{Campbell}-\autor{Mankiw}, 1989). \concepto{Alisamiento}: tasa según $r,\rho,\theta$; nivel según la riqueza total $\Rightarrow$ renta permanente y ciclo vital \vertema{3.A.33}.
\item[] ↑$r$: \emph{sustitución} ↓$C_0$; \emph{renta} ↑$C_0$ si presta; \emph{riqueza} ↓$C_0$. Total ambiguo si presta, negativo si pide. Política: con ESI baja, gravar menos el interés apenas ↑ahorro.
\item[] \emph{Tipo real} (bono 10 años EE.\,UU., IPC): 2,0\,\% (1962-2025); 4,2\,\% (1980-1999); 0,5\,\% (2008-2021); $-0{,}5$\,\% (2022-2025); TIPS 2,60\,\% (30/10/2026).
\item \autor{Becker} (1962): capital humano; \concepto{teorema de separación} (\autor{Fisher}; \autor{Hirshleifer}, 1958; \autor{Acemoglu}-\autor{Autor}, 2011) \vertema{3.A.25}. Con $u$ cuadrática y $\beta(1+r)=1$, \concepto{paseo aleatorio} (\autor{Hall}, 1978); C-CAPM (\autor{Lucas}, 1978; \autor{Breeden}, 1979).
\item \emph{Cota cero}: con efectivo, $i_t\geq0$; si el tipo que vacía exige $i_t<0$, ajustan cantidades (\autor{Krugman}, 1998).
\end{itemize}

\sub{I.2. Decisiones intertemporales de las empresas\label{rep:I.2}}
\begin{itemize}
\item Max. VA de beneficios; si $B_t$ solo depende de $t$, miope. Lo hace intertemporal una ecuación de estado, $K_{t+1}=(1-\delta)K_t+I_t$; sin fricciones, \autor{Jorgenson}: $PF_K=P^K(r+\delta)$, sucesión de estáticos. Hacen falta costes de ajuste o recursos agotables.
\item \textbf{a) Robinson Crusoe}: \emph{sin mercado}, RMT $1+F'(K_{t+1})-\delta$ = RMS: la inversión depende de la impaciencia. \emph{Con mercado}: \autor{Euler} + $F'(K_{t+1})-\delta=r$: \concepto{separación de \autor{Fisher}} (\autor{Hirshleifer}, 1958): primero maximizar valor, luego prestar/pedir; base del VAN (+ \autor{Modigliani}-\autor{Miller}). Falla con restricciones de liquidez.
\item \textbf{b) q de \autor{Tobin}} (1969; \autor{Romer}, cap. 9): costes de ajuste convexos $\Phi(I_t)$, $\Phi(0)=\Phi'(0)=0$:
\begin{formula}
\centering $q_t=1+\Phi'(I_t),\qquad r\,q_t=B(K_t)+\dot q_t$
\end{formula}
\item[] $q$ = precio sombra del capital instalado; se invierte si $q>1$; la 2.ª (derivando respecto al \emph{estado}) = \autor{Euler} de un activo \vertema{3.A.34}.
\item \textbf{c) \autor{Hotelling}} (1931): \concepto{regalía} $p_t-\mathrm{CMg}=\lambda_t$ y
\begin{formula}
\centering $\dfrac{d(p_t-\mathrm{CMg})/dt}{p_t-\mathrm{CMg}}=r$
\end{formula}
\item[] Renta de escasez, no beneficio extraordinario; \emph{paradoja verde} \vertema{3.B.31}.
\end{itemize}

% ==========================================================================
\bloque[8]{II. El agente representativo de vida infinita\label{rep:II}}
\begin{ideaclave}
Llevar las decisiones a la macro exige un agente representativo; el de vida infinita da un equilibrio eficiente que nunca acumula de más: el mercado ahorra lo que conviene.
\end{ideaclave}

\sub{II.1. El agente representativo\label{rep:II.1}}
\begin{itemize}
\item \emph{Por qué}: crítica de \autor{Lucas} (1976): buscar los fundamentos; manejable a costa de la heterogeneidad; base de los EGDE. Historia: \emph{trading bodies} (\autor{Jevons}, 1871), empresa representativa (\autor{Marshall}), crítica de \autor{Robbins} (1928).
\item 3 sentidos: todos idénticos; el «medio»; \emph{el que reproduce las decisiones agregadas} (el útil).
\item SMD: la racionalidad individual apenas restringe la demanda agregada. Hogar representativo \emph{positivo} $\iff$ \concepto{forma polar de \autor{Gorman}} con misma pendiente:
\begin{formula}
\centering $V_i(\vect p,\overline W_i)=a_i(\vect p)+b(\vect p)\,\overline W_i$
\end{formula}
\item[] Engel lineales y paralelas: efectos renta \emph{iguales} (no nulos): homotéticas idénticas o cuasilineales en el mismo bien. \emph{Normativo}: más condiciones.
\item Empresa representativa: fácil (sin efectos renta; sin externalidades, competencia). Sector público: peor (monetaria, regla; fiscal, gasto improductivo).
\item 2 familias: \concepto{MHI} (vida infinita) y \concepto{MGS} (vidas finitas que conviven); \autor{Weil} (1989) combina.
\end{itemize}

\sub{II.2. Los modelos de horizonte infinito\label{rep:II.2}}
\begin{itemize}
\item \emph{Definición}: horizonte no acotado. Natural en la empresa; en el hogar, final lejano o incierto (prob. de muerte constante = inmortal más impaciente) o \concepto{dinastía} altruista (\autor{Barro}, 1974): «de tal palo, tal astilla» $\Rightarrow$ equivalencia ricardiana \vertema{3.A.39}.
\item \emph{Rasgos}: (1) microfundamentados: \emph{menos} expuestos a \autor{Lucas}, no inmunes; (2) 3 condiciones: utilidad finita, \concepto{no-Ponzi} $\lim b_{T+1}/(1+r)^T\geq0$ y \concepto{transversalidad}:
\begin{formula}
\centering $\lim_{T\to\infty}\beta^Tu'(C_T)\,b_{T+1}=0;\quad \lim_{t\to\infty}e^{-\rho t}\lambda_tb_t=0$
\end{formula}
\item[] (necesaria en general: \autor{Michel}, 1982; elige la senda al EE);
\item[] (3) \concepto{consistencia dinámica}: exponencial, el único si el descuento depende solo de la distancia (\autor{Strotz}, 1955). Evidencia: \concepto{hiperbólico} (\autor{Thaler}, 1981; \autor{Loewenstein}-\autor{Prelec}, 1992) o \concepto{cuasihiperbólico} (\autor{Laibson}, 1997) $\Rightarrow$ inconsistencia, infraahorro, mecanismos de compromiso \vertema{3.A.1}.
\end{itemize}
\graficorepaso[0.85\linewidth]{descuento}{Descuento exponencial, hiperbólico y cuasihiperbólico (\autor{Angeletos} et al., 2001): mucho entre hoy y mañana, poco entre 2 fechas lejanas}
\begin{itemize}
\item[] Gobierno con agentes que miran al futuro: inconsistencia de la política (\autor{Kydland}-\autor{Prescott}, 1977) \vertema{3.A.36}.
\item[] (4) \emph{Resolución}: planificador (1TFEB) \vertema{3.A.22}; EE y fases; \concepto{senda de silla}: única solución estable si raíces inestables = variables que saltan (\autor{Blanchard}-\autor{Kahn}, 1980); calibración e impulso-respuesta \vertema{3.A.42}. (5) Análisis de bienestar.
\item \textbf{RCK} (\autor{Ramsey}, 1928; \autor{Cass}, \autor{Koopmans}, 1965): \autor{Solow} con hogares que eligen el ahorro; $L_t=e^{nt}$, $\int e^{-(\rho-n)t}u(C_t/L_t)dt$, utilidad acotada si $\rho-n>(1-\theta)g$; $r_t=f'(k_t)-\delta$:
\begin{formula}
\centering $\dfrac{\dot c_t}{c_t}=\dfrac{f'(k_t)-\delta-\rho-\theta g}{\theta}$\\[2pt]
$\dot k_t=f(k_t)-c_t-(n+g+\delta)k_t$
\end{formula}
\item[] EE único, punto de silla; \concepto{regla de oro modificada} $f'(k^*)=\rho+\delta+\theta g>n+g+\delta=f'(k_{\mathrm{oro}})$: \emph{nunca acumula de más}; equilibrio eficiente \vertema{3.A.43}.
\item \emph{Aplicaciones}: parcial: consumo \vertema{3.A.33}, q \vertema{3.A.34}, dinero (\autor{Sidrauski}, \autor{Clower}, \autor{McCallum}) \vertema{3.A.35}, trabajo \vertema{3.A.25} \vertema{3.A.26}; general: RCK, $AK$, \autor{Uzawa}-\autor{Lucas} \vertema{3.A.44}, ciclo real y EGDE NK \vertema{3.A.7}, deuda, BP intertemporal \vertema{3.B.12}, DICE (\autor{Nordhaus}, 1992).
\end{itemize}

% ==========================================================================
\bloque[10]{III. Generaciones solapadas y heterogeneidad\label{rep:III}}
\begin{ideaclave}
Basta con vidas finitas que se solapan para que el equilibrio competitivo, sin fallo de mercado, pueda acumular demasiado capital; entonces dinero, deuda y reparto pueden mejorar a todas las generaciones.
\end{ideaclave}

\sub{III.1. Los modelos de generaciones solapadas\label{rep:III.1}}
\begin{itemize}
\item \concepto{MGS}: vida finita, convive con los sucesores; la economía no acaba (\autor{Allais}, 1947; \autor{Samuelson}, 1958; \autor{Diamond}, 1965). Rasgos: (1) se solapan hogares, no empresas; (2) equilibrio general (o intercambio puro); (3) heterogeneidad por edad; (4) existencia, unicidad y estabilidad no garantizadas (\autor{Galor}-\autor{Ryder}, 1989); (5) el \emph{1TFEB puede fallar}: \concepto{doble infinitud} (\autor{Shell}, 1971): valor de la dotación agregada infinito; (6) posible ineficiencia dinámica.
\item \textbf{\autor{Samuelson}} (1958), \emph{el dinero como contrato social}: 2 periodos, dotación $e$ perecedera de joven, nada de viejo. Sin dinero, \emph{autarquía} ($C^J=e$, $C^V=0$). Factible en EE:
\begin{formula}
\centering $C^J+\dfrac{C^V}{1+n}=e$
\end{formula}
\item[] \emph{tipo de interés biológico} $n$. Con $M$ fijo, $P_{t+1}/P_t=1/(1+n)$, el dinero rinde $1+n$: $u'(C^J)=\beta(1+n)u'(C^V)$; con $\ln$, $C^J=e/(1+\beta)$. Autarquía = equilibrio no Pareto; el monetario mejora a todos. Valor del dinero = burbuja de confianza; reparto, igual: «cuando seas padre, comerás huevos». Límites: no medio de cambio, equilibrios múltiples, sin capital.
\end{itemize}
\graficorepaso[0.6\linewidth]{samuelson}{\autor{Samuelson}: autarquía $A$ y equilibrio monetario $M$ (pendiente $-(1+n)$)}
\begin{itemize}
\item \textbf{\autor{Diamond}} (1965): \autor{Samuelson} + producción de \autor{Solow}. Jóvenes trabajan, consumen y ahorran $S_t$; viejos, $C^V_{t+1}=(1+r_{t+1})S_t$; $\ln C^J+\beta\ln C^V$; $f=k^\alpha$, $\delta=1$ (periodos de 30 años):
\begin{formula}
\centering $S_t=\dfrac{\beta}{1+\beta}\,w_t$ (no depende de $r$)
\end{formula}
\item[] $K_{t+1}=L_tS_t$ $\Rightarrow$ \concepto{ley de movimiento}:
\funcionclave{k_{t+1}=\frac{\beta(1-\alpha)}{(1+\beta)(1+n)(1+g)}\,k_t^{\alpha}}
\item[] Cóncava, único EE estable (pendiente $\alpha<1$); ↑$\beta$ ↑$k^*$; ↑$n$, $g$ ↓$k^*$. Ahorro $=\beta(1-\alpha)/(1+\beta)$. CRRA general: $S$ ↑ con $r$ si $\theta<1$.
\end{itemize}
\graficorepaso[0.8\linewidth]{diamond}{\autor{Diamond}: ley de movimiento cóncava, bisectriz y EE estable; ineficiencia si $k^*>k_{\mathrm{oro}}$}
\begin{itemize}
\item \concepto{Regla de oro} (\autor{Phelps}, 1961): $f'(k_{\mathrm{oro}})-\delta=(1+n)(1+g)-1\simeq n+g$. \concepto{Ineficiencia dinámica} si
\begin{formula}
\centering $r^*<(1+n)(1+g)-1\simeq n+g$
\end{formula}
\item[] Reducir $k$ ↑ el consumo de \emph{todas} las generaciones: «la avaricia rompe el saco». Con $\delta=1$: $\beta>\alpha/(1-2\alpha)$; con $\alpha=1/3$, nunca; con 0,3, basta $\beta>0{,}75$ ($\rho<1$\,\% anual). Causa: ahorrar para la vejez solo con capital; no se presta a los no nacidos.
\item \emph{Extensiones}: \emph{deuda} desplaza capital: ↓ bienestar si eficiente, ↑ si ineficiente (\autor{Blanchard}, 2019) \vertema{3.A.39}; \emph{equivalencia ricardiana} falla salvo herencias altruistas (\autor{Barro}, 1974); \emph{pensiones}: capitalización rinde $1+r$ (capital igual); \emph{reparto} rinde $(1+n)(1+g)$ y ↓ capital; mejor si $n+g>r$: \concepto{condición de \autor{Aaron}} (1966) = ineficiencia dinámica \vertema{3.A.38} \vertema{4.B.24}; \emph{burbujas} (\autor{Tirole}, 1985): solo si sin ella es ineficiente; la elimina; \autor{Blanchard} (1985)-\autor{Yaari} (1965): prob. de muerte constante, deuda no neutral; \autor{Auerbach}-\autor{Kotlikoff} (1987).
\item \emph{Envejecimiento} ↓$n$ ↓rendimiento del reparto: dependencia España 31 (2025) $\to$ 55 (2050, Eurostat; INE 52); gasto en pensiones 13,1\,\% (2022) $\to$ 17,3\,\% (2050) vs. 12,6\,\% zona del euro.
\item \emph{Evidencia}: $r-g$ (deuda) $<0$ casi siempre desde 2000: España $-0{,}8$ (2000-2025), zona del euro $-0{,}5$, EE.\,UU. $-0{,}8$ (desde 1962); $>0$ en 1980-1999 (+1,8 EE.\,UU., +1,7 España) y en la crisis de la deuda.
\end{itemize}
\graficorepaso[0.9\linewidth]{r-menos-g}{$r-g$ (tipo real a largo menos crecimiento), medias de 5 años: España, zona del euro y EE.\,UU. (AMECO)}
\begin{itemize}
\item[] Pero $r<g$ en la deuda \emph{no prueba} ineficiencia: cuenta $f'(k)-\delta$. \concepto{Criterio de \autor{Abel}-\autor{Mankiw}-\autor{Summers}-\autor{Zeckhauser}} (1989): rentas del capital $>$ inversión $\Rightarrow$ eficiente; 2025: rentas 43-46\,\% del PIB, inversión 21-22\,\%. \autor{Geerolf} (2018): Japón y Corea, quizá de más.
\item \emph{Cuadro MHI vs. MGS}: hogar inmortal vs. 2 periodos; \autor{Euler} vs. $S=\beta w/(1+\beta)$; siempre eficiente vs. no si $r^*<n+g$; 1TFEB sí vs. puede fallar; ricardiana sí vs. no; dinero/deuda/reparto neutrales vs. pueden mejorar; ciclo y crecimiento vs. demografía, pensiones, deuda.
\end{itemize}

\sub{III.2. Limitaciones y nuevos enfoques\label{rep:III.2}}
\begin{itemize}
\item Validez teórica limitada (\autor{Gorman}, SMD; analogía con \autor{Arrow} \vertema{3.A.24}); puede engañar: reacciona al revés que todos (\autor{Kirman}, 1992).
\item \emph{Nuevos}: heterogéneos con mercados incompletos (\autor{Aiyagari}, 1994; \autor{Krusell}-\autor{Smith}, 1998); \concepto{HANK} (\autor{Kaplan}-\autor{Moll}-\autor{Violante}, 2018): la política monetaria actúa por la renta de los hogares con poca liquidez, no por \autor{Euler}; 2 tipos de hogares (\autor{Galí}-\autor{López-Salido}-\autor{Vallés}, 2007) \vertema{3.A.37}; basados en agentes; conocimiento imperfecto (\autor{Frydman}-\autor{Goldberg}); comportamiento (\autor{De Grauwe}); \autor{Wren-Lewis} (2018): convivencia.
\end{itemize}

% ==========================================================================
\bloque[1]{Conclusión}
\begin{itemize}
\item \textbf{Recapitulación}: «En conclusión,» variable de estado $\Rightarrow$ dinámica; 3 métodos, una \autor{Euler} y un precio sombra; MHI eficiente, nunca de más; MGS puede ahorrar de más y dinero, deuda y reparto mejorar a todos.
\item \textbf{Relevancia}: cimientos de la macro actual (bancos centrales, sostenibilidad de deuda y pensiones, clima).
\item \textbf{Extensiones}: \vertema{3.A.33} \vertema{3.A.34} \vertema{3.A.43} \vertema{3.A.44} \vertema{3.A.42} \vertema{3.A.39} \vertema{3.A.36} \vertema{3.A.14} \vertema{3.A.1} \vertema{3.B.12}.
\item \textbf{Opinión}: la macro debe más a \autor{Fisher} y \autor{Ramsey} que a \autor{Keynes} (matizado por las rigideces NK); las preguntas urgentes (pensiones, deuda, clima) son entre generaciones: hacen falta modelos con generaciones.
\item \textbf{Idea final}: dejar para mañana lo que mañana valga más (= \autor{Euler}); ¿cuánto ahorrar? Depende de cómo modelicemos el tiempo: generaciones que se suceden pueden ahorrar de más, y entonces es avaricia que rompe el saco.
\end{itemize}

% ==========================================================================
\bloque{Autores y fechas}
\hito{1871}{\autor{Jevons}: \emph{trading bodies}}
\hito{1928}{\autor{Ramsey}: ahorro óptimo; \autor{Robbins}: crítica de la empresa representativa}
\hito{1930}{\autor{Fisher}: teoría del interés, separación}
\hito{1931}{\autor{Hotelling}: recursos agotables}
\hito{1937}{\autor{Samuelson}: descuento exponencial}
\hito{1947}{\autor{Samuelson}: principio de correspondencia; \autor{Allais}}
\hito{1953}{\autor{Gorman}: forma polar}
\hito{1955}{\autor{Strotz}: consistencia temporal}
\hito{1957}{\autor{Bellman}: programación dinámica}
\hito{1958}{\autor{Samuelson}: MGS y dinero; \autor{Hirshleifer}}
\hito{1961}{\autor{Phelps}: regla de oro}
\hito{1962}{\autor{Pontryagin} et al.: principio del máximo; \autor{Becker}}
\hito{1965}{\autor{Cass}, \autor{Koopmans}; \autor{Diamond}; \autor{Yaari}}
\hito{1966}{\autor{Aaron}: condición del reparto}
\hito{1969}{\autor{Tobin}: q}
\hito{1971}{\autor{Shell}: doble infinitud}
\hito{1974}{\autor{Barro}: dinastía y deuda}
\hito{1976}{\autor{Lucas}: crítica}
\hito{1977}{\autor{Kydland}-\autor{Prescott}: inconsistencia}
\hito{1978}{\autor{Hall}: paseo aleatorio}
\hito{1980}{\autor{Blanchard}-\autor{Kahn}}
\hito{1985}{\autor{Tirole}: burbujas; \autor{Blanchard}: juventud perpetua}
\hito{1989}{\autor{Abel} et al.: criterio de eficiencia}
\hito{1992}{\autor{Kirman}; \autor{Nordhaus}: DICE}
\hito{1997}{\autor{Laibson}: cuasihiperbólico}
\hito{2018}{\autor{Kaplan}-\autor{Moll}-\autor{Violante}: HANK; \autor{Geerolf}}

% ==========================================================================
\bloque{Datos clave}
\begin{itemize}
\item Mecanismo de equidad intergeneracional: \dato{0,90 p.p. (2026) $\to$ 1,2 (2029-2050)}{LGSS}.
\item Tipo real 10 años EE.\,UU.: \dato{2,0\,\% (1962-2025)}{FRED}; TIPS \dato{2,60\,\%}{Tesoro, 30/10/2026}.
\item $r-g$ medio: España \dato{$-0{,}8$ (2000-2025)}{AMECO}; zona del euro \dato{$-0{,}5$}{AMECO}; 2025, España \dato{$-2{,}5$}{AMECO}.
\item Rentas del capital \dato{43-46\,\% del PIB}{AMECO, 2025} vs. inversión \dato{21-22\,\%}{AMECO}.
\item Dependencia España: \dato{31 (2025) $\to$ 55 (2050)}{Eurostat}.
\item Gasto en pensiones España: \dato{13,1\,\% (2022) $\to$ 17,3\,\% (2050)}{CE, \emph{Ageing Report} 2024}.
\end{itemize}

% ==========================================================================
\bloque{Preguntas del tribunal}
\begin{enumerate}
\item \textbf{¿Cuándo es genuinamente dinámico un problema?} Cuando una variable de estado une los periodos; si no, es una sucesión de estáticos (la empresa sin costes de ajuste).
\item \textbf{¿Por qué los 3 métodos dan lo mismo?} Resuelven el mismo problema; $\lambda$, $\lambda_t$ y $V'(b_t)$ son el mismo precio sombra; cambia la comodidad.
\item \textbf{¿Para qué la transversalidad?} Elige, entre las sendas que cumplen el resto, la que va al EE; no dejar nada de valor al final. No-Ponzi: no endeudarse sin límite.
\item \textbf{¿Por qué RCK no es ineficiente?} Utilidad acotada $\Rightarrow$ $\rho+\theta g>n+g$, luego $k^*<k_{\mathrm{oro}}$.
\item \textbf{¿Por qué puede serlo \autor{Diamond} sin fallo de mercado?} Solo se ahorra para la vejez con capital y no se presta a los no nacidos; doble infinitud.
\item \textbf{¿$r<g$ prueba ineficiencia?} No: cuenta la rentabilidad del capital; criterio de \autor{Abel} et al.: rentas $>$ inversión.
\end{enumerate}

\begin{trampa}
\begin{itemize}
\item ESI $=1/\theta$; $\theta$ también es aversión relativa al riesgo; $\theta=1$, logarítmica (no «máxima concavidad»).
\item Con $\ln$: ahorro de \autor{Diamond} y fracción consumida en \autor{Bellman} \emph{no} dependen de $r$.
\item Alisamiento: se da con cualesquiera $r$ y $\rho$ (examen de julio de 2021, pregunta 7, a)); $r>\rho$ no implica ser prestamista.
\item \autor{Gorman}: efectos renta \emph{iguales}, no nulos.
\item Exponencial = único consistente; en el cuasihiperbólico $\beta$ es sesgo al presente.
\item q: $q=1+\Phi'(I_t)$ (coste \emph{marginal}); la condición dinámica, respecto al estado.
\item \autor{Hotelling}: crece la \emph{regalía} (precio neto), no el precio, a la tasa $r$.
\item Ineficiencia: $r^*<n+g$ (no solo $n$); \autor{Aaron}: reparto mejor si $n+g>r$.
\item RCK: regla de oro \emph{modificada} $\rho+\delta+\theta g$; \autor{Solow}: $n+g+\delta$.
\item MHI: menos expuestos a \autor{Lucas}, no inmunes.
\item Principio del máximo: anunciado en 1956, publicado en 1962.
\end{itemize}
\end{trampa}

\end{columnas}
\end{document}
